\begin{figure}[H]
\centering
\includegraphics[width=1.0\linewidth]{wakayama/A2_landing_map.pdf}
\caption[落下点を地図に重ねる]{★の「和歌山中央」（\ang{33.902}N, \ang{135.375}E）から
冬に上げたスペースバルーンの落下点を、
標高と水深に重ねたもの。各条件 \num{4000} 点、観測は \num{3612} 点。
条件の中身は表 \ref{tab:winter-scenarios} にある。
点の描画順は広い雲が先、狭い雲が後なので、
いちばん狭いマゼンタ（層独立）が最前面にある。
地形・水深は NOAA NCEI の DEM を陸と海で別の色階調にし、陰影を重ねたもの。
海岸線は標高 0 の等値線である。}
\label{fig:a2-wakayama}
\end{figure}
